\begin{textsample}{POS Dim 3 – pr – Score 15.00 – pr/pr\_inf\_1.txt}  \label{ex:f3_pos_253}
Educação \textbf{Infantil} A partir da homologação da Base Nacional Comum Curricular ( BNCC ) para a Educação \textbf{Infantil}, apresenta -se o desafio da elaboração de um documento de orientação às instituições de ensino que ofertam essa etapa da Educação Básica, incorporando as determinações legais do documento normativo e respeitando as características do território paranaense.
Concorrem para esta elaboração se constituir em desafio fatores como a diversidade sociocultural e estrutural dos municípios, das redes e dos sistemas de ensino e seus diferentes projetos para atendimento das \textbf{crianças} de 0 a 5 anos nos Centros de Educação \textbf{Infantil} e escolas.
No entanto, um dos fatores comuns a todos é o compromisso de atender, com qualidade, a ampliação da oferta da Educação \textbf{Infantil} instituída na Meta 1 do Plano Estadual de Educação do Paraná. ( PARANÁ, 2015, p. 58-59 ).
Nesse sentido, a BNCC avança como elemento de interlocução entre as redes municipais, a rede estadual e as redes privadas que buscam a melhoria da qualidade na Educação \textbf{Infantil}, promovendo a equidade das práticas pedagógicas \textbf{apoiadas} nos direitos e objetivos de aprendizagem e desenvolvimento.
Para isso, é importante que os sistemas e redes de ensino e as instituições escolares organizem seus planejamentos com foco no reconhecimento das necessidades dos estudantes, levando em conta suas diferenças e priorizando o acesso, permanência e sucesso de todos os alunos, independentemente de sua condição.
Um dos indicadores de qualidade é a existência, em cada instituição, de um projeto político pedagógico elaborado e revisado constantemente pelos profissionais que nela atuam, considerando “ as orientações legais vigentes e [... ] os conhecimentos já acumulados a respeito da educação \textbf{infantil} ” ( BRASIL, 2009, p. 37 ).
É no Projeto Político Pedagógico que se consolida o currículo e se definem as especificidades para o trabalho articulado entre o cuidar e o educar inerente à Educação Básica.
Os Projetos Políticos Pedagógicos da Educação \textbf{Infantil}, desde 2009, orientados pelas Diretrizes Curriculares Nacionais para a Educação \textbf{Infantil} ( DCNEIs ), estabelecidas pela Resolução nº 5/2009 – CNE/CEB, devem ter como seus eixos norteadores as interações e a \textbf{brincadeira}.
Essa orientação é confirmada na BNCC, pois são “ experiências nas quais as \textbf{crianças} podem construir e apropriar -se de conhecimentos por meio de suas ações e interações com seus pares e com os \textbf{adultos}, o que possibilita aprendizagens, desenvolvimento e socialização ” ( BRASIL, 2017, p. 35 ). \textbf{Brincadeiras} e interações acontecem diariamente entre as \textbf{crianças} e representam o direito à infância, a viver e crescer em um ambiente \textbf{lúdico} e prazeroso que lhes proporcione segurança e confiança.
Mas isso não significa que esses momentos dispensem a necessidade de intencionalidade e planejamento da prática pedagógica, pois os objetivos de aprendizagem e desenvolvimento se tornam mais complexos ou diferentes em cada faixa etária.
Nesse sentido, é importante planejar considerando as singularidades e o direito de aprender de todos.
Além dos eixos interações e a \textbf{brincadeira}, a BNCC, compreendendo a \textbf{criança} por inteiro – corpo, mente e emoções, aponta a importância de conviver, \textbf{brincar}, participar, explorar, expressar e conhecer -se como direitos essenciais de aprendizagem e desenvolvimento.
A estruturação dos currículos prevista na BNCC com uma organização em campos de experiência reafirma as DCNEIs, em especial o seu artigo 3º : > O currículo da Educação \textbf{Infantil} é concebido como um conjunto de práticas que buscam articular as experiências e os saberes das \textbf{crianças} com os conhecimentos que fazem parte do patrimônio cultural, artístico, ambiental, científico e tecnológico, de modo a promover o desenvolvimento integral de \textbf{crianças} de 0 a 5 anos de idade. ( BRASIL, 2009, p. 1 ).
Os objetivos de aprendizagem estão organizados em cinco campos de experiências : O eu, o outro e o nós ; Corpo, \textbf{gestos} e movimentos ; Traços, \textbf{sons}, cores e formas ; \textbf{Escuta}, fala, pensamento e imaginação ; e Espaços, tempos, quantidades, relações e transformações.
Esses campos “ constituem um arranjo curricular que \textbf{acolhe} as situações e as experiências concretas da vida cotidiana das \textbf{crianças} e seus saberes, entrelaçando -os aos conhecimentos que fazem parte do patrimônio cultural ” ( BRASIL, 2017, p. 38 ).
Essa é uma forma de fortalecer a Educação \textbf{Infantil} com a sua especificidade no trabalho educativo, não confundindo com práticas antecipatórias e preparatórias, que pouco contribuem para o processo formativo da \textbf{criança}.
Na sua estruturação, a BNCC define agrupamentos para as \textbf{crianças} em três fases, sendo estas : \textbf{bebês}, \textbf{crianças} bem \textbf{pequenas} e \textbf{crianças} \textbf{pequenas}.
Essas fases consideram a proximidade dos objetivos, “ que correspondem, aproximadamente, às possibilidades de aprendizagem e às características do desenvolvimento das \textbf{crianças} ” ( BRASIL, 2017, p. 42 ).
Partindo das premissas da BNCC sobre os eixos integradores, os direitos essenciais de aprendizagem e desenvolvimento, os campos de experiências e a estruturação dos currículos, se pauta a elaboração do Referencial Curricular do Paraná : princípios, direitos e orientações.
O documento está organizado em cinco seções, nas quais se apresentam : Considerações históricas da Educação \textbf{Infantil}, Princípios básicos da Educação \textbf{Infantil} e os direitos de aprendizagem, Concepções norteadoras do trabalho pedagógico na Educação \textbf{Infantil}, Articulação entre Educação \textbf{Infantil} e Ensino Fundamental e Organizador Curricular.
Desta forma, em função da relevância de contextualizar a Educação \textbf{Infantil}, inicia -se com uma breve retomada da sua história no Brasil, desde os Jardins da Infância iniciados ao final do século XIX, até sua inserção na Educação Básica.
Observa -se que seu reconhecimento enquanto direito das \textbf{crianças} e das famílias é o resultado de intensas lutas sociais pela sua inclusão nas legislações e nas políticas públicas brasileiras.
Em seguida, são apresentados os princípios básicos da Educação \textbf{Infantil} segundo as DCNEIs, \textbf{demonstrando} a articulação destes com os direitos de aprendizagem estabelecidos na BNCC.
A concepção de \textbf{criança} que pauta o documento é apresentada na sequência.
Os tópicos seguintes seguem a estrutura da BNCC apresentada anteriormente, trazendo os eixos interações e a \textbf{brincadeira} inseridos nos campos de experiências.
A necessária articulação entre as etapas Educação \textbf{Infantil} e Ensino Fundamental é abordada buscando incentivar as redes e instituições quanto ao planejamento de ações para tornar essa transição adequada, garantindo o direito de infância.
Isso significa que é preciso evitar rupturas nesse processo, privilegiando a relação e diálogo entre as etapas envolvidas, tanto na elaboração dos currículos, como em sua prática.
O organizador curricular é apresentado na sequência, trazendo os objetivos estabelecidos pela BNCC enriquecidos com desdobramentos, tornando -os mais específicos para cada idade que compõe os agrupamentos.
A opção de estabelecer objetivos por idade, e não pelos agrupamentos indicados na BNCC, está pautada na autonomia das redes e sistemas de ensino quanto à formação das turmas.
Outra característica presente no organizador curricular é a inclusão de “ saberes e conhecimentos ” como elementos que, associados aos campos, aglutinam uma série de objetivos próximos e marcam a intencionalidade das práticas docentes que oportunizam a construção de “ sentidos sobre a natureza e a sociedade, produzindo cultura ” ( BRASIL, 2009, p. 1 ).
Finalmente, é importante descrever o caminho percorrido e a sistemática da elaboração do Referencial Curricular Paranaense para a etapa da Educação \textbf{Infantil}.
Por determinação do Ministério da Educação ( MEC ) através da Portaria nº 331/2018 – SEB/MEC, cabe ao Estado, em regime de colaboração com os Municípios, a implementação da BNCC no território paranaense segundo os critérios estabelecidos na mesma.
Assim, constituiu -se o Comitê Executivo Estadual e a Assessoria Técnica Estadual para implementação da BNCC, cujas funções foram, dentre outras, indicar e orientar um grupo de especialistas em Educação \textbf{Infantil} para redigir a versão preliminar.
Esta versão, posta em consulta pública e também objeto de estudo de eventos de formação continuada das redes de ensino, recebeu contribuições dos educadores e da comunidade se tornando um documento representativo das aspirações e necessidades da Educação \textbf{Infantil} no Estado do Paraná, inclusive com oportunidade de rever as questões relacionadas à inclusão, igualdade, diversidade e equidade, as quais devem permear todo o trabalho.
Toda prática pedagógica desenvolvida na Educação \textbf{Infantil} deve considerar a identificação das necessidades educacionais especiais dos estudantes, bem como as intervenções necessárias em ambiente de ensino para o atendimento de todos os alunos.
Dessa maneira, a BNCC e o Referencial Curricular do Paraná, além de trazerem a obrigatoriedade da elaboração ou reorganização curricular, recolocam na pauta das políticas públicas a discussão sobre a infância e sobre a necessidade de aprofundamento dos fundamentos e concepções que amparam as práticas pedagógicas na Educação \textbf{Infantil}.

% matched lemmas: acolher, adulto, apoiar, bebê, brincadeira, brincar, criança, demonstrar, escuta, gesto, infantil, lúdico, pequeno, som
\end{textsample}
